\textbf{UC-ROZ-03: Kompensata i saldo końcowe}

\begin{longtable}{|p{4cm}|p{11cm}|}
\hline
\textbf{Element} & \textbf{Opis} \\
\hline

Identyfikator & UC-ROZ-03 \\
\hline

Nazwa & Kompensata i saldo końcowe \\
\hline

Opis &
Proces agregacji wpłat klienta, finansowania bankowego oraz wartości odkupu pojazdu w celu wyliczenia końcowej kwoty wymaganej do zapłaty przed wydaniem pojazdu. \\
\hline

Aktor główny & Księgowy / Handlowiec \\
\hline

Aktorzy pomocniczy & Kontekst Sprzedaży i CRM, Kontekst Finansowania i ubezpieczeń \\
\hline

Warunki wstępne &
Pojazd posiada status gotowy do wydania, a wszystkie wcześniejsze operacje finansowe zostały zarejestrowane w systemie. \\
\hline

Warunki końcowe &
System generuje końcowe saldo rozliczenia oraz dokument rozliczeniowy dla klienta. \\
\hline

Scenariusz główny &
\begin{enumerate}
\item Handlowiec inicjuje proces rozliczenia końcowego przed wydaniem pojazdu.
\item System pobiera całkowitą wartość brutto pojazdu z zamówienia.
\item System pobiera i agreguje wszystkie wcześniej zaksięgowane wpłaty klienta.
\item System komunikuje się z Kontekstem Finansowania i ubezpieczeń, pobiera wartość zatwierdzonego finansowania bankowego.
\item System weryfikuje istnienie odkupu pojazdu używanego oraz pobiera jego wartość rozliczeniową.
\item System oblicza saldo końcowe według reguły biznesowej:
\[
\text{Saldo końcowe} =
\text{Wartość pojazdu}
- \text{Zadatki}
- \text{Finansowanie}
- \text{Wartość odkupu}
\]
\item System generuje dokument rozliczeniowy zawierający końcową kwotę do zapłaty przez klienta.
\end{enumerate}
\\
\hline

Scenariusze alternatywne &
\begin{itemize}
\item [A1] Nadpłata: W kroku 6 system wylicza saldo ujemne. System blokuje standardowe rozliczenie, oznacza zamówienie statusem „Wymaga korekty księgowej” oraz generuje dyspozycję zwrotu nadpłaty.
\item [A2] Brak potwierdzenia finansowania: W kroku 4 system nie otrzymuje potwierdzenia uruchomienia środków leasingowych lub kredytowych. System blokuje możliwość wydania pojazdu do czasu uzyskania statusu „Finansowanie zatwierdzone”.
\end{itemize}
\\
\hline

\end{longtable}